\documentclass[../main.tex]{subfiles}
\graphicspath{{\subfix{../image/}}}
\begin{document}
\chapter{Lie 群}\label{chapter:ch07-Lie群}

本章引入 Lie 群，它们既是光滑流形，又是群，且乘法与取逆都是光滑映射。Lie 群本身提供了许多有趣流形的例子，同时也是研究更一般流形的基本工具，这主要是因为它们在作为其他流形的对称群时所起的作用。

本章的目的是介绍 Lie 群以及处理它们的一些工具，并给出丰富的例子。在后续章节中（尤其是第 8、20、21 章），我们将发展 Lie 群的更多性质与应用。

我们从 Lie 群的定义以及与之相关的一些基本结构开始，然后给出若干例子。接下来研究 Lie 群同态，即同时也是光滑映射的群同态。然后引入 Lie 子群（同时也是光滑子流形的子群），由此产生许多新的 Lie 群例子。

在解释这些基本概念之后，我们引入 Lie 群在流形上的作用，这是 Lie 群存在的主要理由。在章末，我们简要提及群表示。

Lie 群的研究由挪威数学家 Sophus Lie 于十九世纪末开创。受法国代数学家 Évariste Galois 发明群论并用其分析多项式方程这一方式的启发，Lie 希望利用以群作用形式表达的对称性来简化偏微分方程与几何中的问题。然而，Lie 当时不可能设想出我们如今称为 Lie 群的整体对象，原因很简单：诸如流形（甚至拓扑空间！）这样的整体拓扑概念尚未形成。Lie 所研究的本质上是 Lie 群的局部坐标版本，如今称为局部 Lie 群。尽管受到他所处时代的局限，他仍为我们如今对 Lie 群的理解奠定了大量基础。我们将在第 20 章中描述他的主要结果（见\cref{theorem:ch20-16}）。

\subfile{Chapter-07/section-01.tex}
\subfile{Chapter-07/section-02.tex}
\subfile{Chapter-07/section-03.tex}
\subfile{Chapter-07/section-04.tex}
\subfile{Chapter-07/section-05.tex}
\end{document}